\section{Introduction}
More than half of the world's population lives in urban areas, where dense materials, reduced vegetation, and anthropogenic heat can amplify exposure during hot weather. The physical basis of the urban heat island has long been understood as an imbalance in radiation, heat storage, and turbulent exchange \cite{Oke1982}. At neighborhood scales, however, the thermal conditions a pedestrian experiences can vary sharply over tens of metres because buildings and vegetation alter short-wave and long-wave radiation.

Urban trees are a prominent heat-adaptation strategy. Their crowns intercept solar radiation and their leaves support evapotranspiration, producing measurable cooling at both local and neighborhood scales \cite{Bowler2010,Armson2012}. The relationship is not spatially uniform: interactions among tree canopy, impervious surface, and the scale of analysis can alter the apparent magnitude of cooling \cite{Ziter2019}. A citywide percentage therefore describes the amount of canopy but not necessarily the shade encountered along an individual route.

Mobile measurement studies can capture human-scale variation that fixed weather stations miss. Yet route-average analyses commonly compress tree cover into a single percentage. Two routes with identical total cover may differ materially: one may provide several long shaded segments, while another alternates rapidly between small crown patches and exposed gaps. We call this route property \emph{canopy continuity}. It is conceptually distinct from total cover and is relevant to pedestrians because short exposed gaps may produce rapid gains in radiant heat load even when route-average air temperature changes little.

This study asks whether canopy continuity explains synthetic pedestrian heat exposure after total canopy cover and urban form are considered. We designed a transparent simulation to demonstrate a defensible mobile-sensor workflow without presenting fabricated observations as empirical evidence. We hypothesized that (1) greater canopy cover would be associated with lower route-average wet-bulb globe temperature, (2) longer uninterrupted canopy segments would provide an additional reduction, and (3) both associations would be stronger during afternoon periods. The synthetic results supported all three hypotheses.
